\SetPageTheme{story}
\PageHeader{Lectia 8}{Verificarea palindromului si exercitii apropiate}{L8 • P1}

\begin{missionbox}[Legam totul]
  Acum avem toate piesele: stim sa extragem cifre, stim sa repetam, stim sa construim oglinditul. Mai ramane doar sa pastram numarul initial si sa facem comparatia finala.
\end{missionbox}

\begin{conceptbox}[Algoritmul de palindrom]
\begin{lstlisting}[style=codequest]
int n, copie, uc, ogl = 0;
citeste n;
copie = n;

cat timp (n != 0) executa
    uc = n % 10;
    ogl = ogl * 10 + uc;
    n = n / 10;
sfarsit cat timp

daca (ogl == copie) atunci
    afiseaza "DA";
altfel
    afiseaza "NU";
sfarsit daca
\end{lstlisting}
\end{conceptbox}

\begin{guidebox}[Prima intrebare]
  De ce avem nevoie de variabila \texttt{copie}?
  \AnswerSpace{2}
\end{guidebox}

\newpage
\SetPageTheme{guided}
\PageHeader{Lectia 8}{Intelegem comparatia finala}{L8 • P2}

\begin{examplebox}[Doua exemple]
  \texttt{n = 1221} produce \texttt{ogl = 1221}, deci raspunsul este \texttt{DA}.

  \texttt{n = 1231} produce \texttt{ogl = 1321}, deci raspunsul este \texttt{NU}.
\end{examplebox}

\begin{guidebox}[Exercitii ghidate]
  1. Ce valoare va avea \texttt{ogl} pentru \texttt{n = 404}?

  2. Ce raspuns se va afisa pentru \texttt{n = 909}?

  3. Ce raspuns se va afisa pentru \texttt{n = 908}?
  \AnswerSpace{5}
\end{guidebox}

\newpage
\SetPageTheme{guided}
\PageHeader{Lectia 8}{Exercitii apropiate algoritmului}{L8 • P3}

\begin{solobox}[Semi-ghidat]
  1. Scrie un algoritm care calculeaza suma cifrelor.

  2. Scrie un algoritm care determina cifra maxima.

  3. Scrie un algoritm care numara cate cifre pare are un numar.
  \AnswerSpace{8}
\end{solobox}

\begin{guidebox}[Ideea comuna]
  In toate aceste probleme scheletul buclei ramane acelasi. Se schimba doar ceea ce facem cu \texttt{uc}.
\end{guidebox}

\newpage
\SetPageTheme{challenge}
\PageHeader{Lectia 8}{Independent}{L8 • P4}

\begin{solobox}[Fisa independenta]
  1. Scrie un algoritm care afiseaza toate numerele palindrom dintre 100 si 200.

  2. Scrie un algoritm care spune daca un numar contine cifra 7.

  3. Explica ce parti din schema de palindrom raman identice si ce parti se schimba.
  \AnswerSpace{9}
\end{solobox}
